\subsection{From an arbitrary decomposition to canonical component bounds}
\label{sec:canonical-cost}
We use the compensator theorem in the following precise form.
For an adapted zero-start c\`adl\`ag process $D$ of integrable total variation,
constant after a finite deterministic time, there is a predictable zero-start
c\`adl\`ag process $D^p$ of integrable variation such that $D-D^p$ is a
uniformly integrable martingale. It is the \emph{dual predictable projection},
characterized by
\[
 E\int K\,dD^p=E\int K\,dD
 \quad\text{for every bounded predictable }K.
\]
This is the compensator theorem, not the assertion that
$D^p_t=E[D_t\mid\mathcal F_{t-}]$.
Existence is the standard dual-projection theorem; the characterization and
jump formula are recalled in \cite[Section 4.3]{Nikeghbali}.
Its variation bound is
\begin{equation}\label{eq:compensator-variation}
 E\Var(D^p)\leq E\Var(D).
\end{equation}
Indeed, apply the defining identity to a predictable Hahn sign for $dD^p$;
the left side is its expected total variation and the right side is bounded
by $E\Var(D)$. Equivalently compensate the two increasing Jordan components
separately and use positivity. For $K=1_F1_{(s,t]}$, $F\in\mathcal F_s$,
the same identity proves the martingale property of $D-D^p$.
The Lean construction obtains this compensator by truncation and limits of
predictable finite-variation approximants. In this account the existence
theorem is used explicitly as a standard foundational result.

\begin{lemma}[Canonical cost of a constant-after-horizon gain]
\label{lem:canonical-cost}
Let $X$ be zero-start and c\`adl\`ag, constant after $T<\infty$, with a
canonical decomposition $X=M^{\rm can}+A^{\rm can}$.
Under the usual conditions,
\[
 \begin{aligned}
 E\Var_{[0,T]}(A^{\rm can})&\leq12\mathcal J(X),\\
 E\sup_{t\leq T}|M^{\rm can}_t-M^{\rm can}_0|
       +E\Var_{[0,T]}(A^{\rm can})&\leq30\mathcal J(X).
 \end{aligned}
\]
These are extended-value inequalities. The infimum defining $\mathcal J$
allows adapted, rather than necessarily predictable, FV components.
\end{lemma}
\begin{proof}
Fix a finite-cost decomposition in that infimum, of cost $j$.
The strict-prefix witness construction gives, at the deterministic stop
$\sigma=T+1$, components $N,B$ representing $X$ before $\sigma$, with
\[
 c_N=E\sup_t|N^\sigma_t|,\qquad
 c_B=E\Var(B^{\sigma-}),\qquad c_N+c_B\leq6j.
\]
Keep their initial values. With $J^N_t=\Delta N_\sigma1_{[\sigma,\infty)}(t)$,
put
\[
 D=B^{\sigma-}-B_0-J^N.
\]
This is adapted and of integrable variation, since
\[
 E\Var(D)\leq c_B+E|\Delta N_\sigma|\leq c_B+2c_N.
\]
Compensate this specific remainder. The processes
\[
 \widehat M=N^\sigma+D-D^p,\qquad \widehat A=B_0+D^p
\]
form a special decomposition of $(N+B)^{\sigma-}=X$.
The last equality uses constancy of $X$ after $T$, not just agreement on
an open interval. By \eqref{eq:compensator-variation},
\[
 E\Var(\widehat A)\leq c_B+2c_N\leq2(c_N+c_B)\leq12j.
\]
Now, and only now, canonical uniqueness identifies the centered components
of this constructed decomposition with those of $X$.
Take the infimum over $j$ to obtain the first bound.
Zero start gives
\[
 \sup_{t\leq T}|M^{\rm can}_t-M^{\rm can}_0|
 \leq\sup_t|X_t|+\Var_{[0,T]}(A^{\rm can}).
\]
The already proved maximal estimate $E\sup_t|X_t|\leq6\mathcal J(X)$
therefore gives $6+12+12=30$. Infinite costs cause no difficulty.
\end{proof}

For a prelocal witness for a regular zero-start $X$ before $\tau\leq T$,
let $c$ be its sum of expected martingale-maximum and variation costs.
The witness comparison in the $r^1$ localization subsection below gives
$\mathcal J_{\tau-}(X)\leq16c$.
The boundary-jump correction in the quadratic-cost subsection gives
\[
 \mathcal J(X^\tau)\leq2\mathcal J_{\tau-}(X)+E|\Delta X_\tau|
 \leq32c+E|\Delta X_\tau|.
\]
If an independently constructed actual integral has gain $X^\tau$, apply
Lemma~\ref{lem:canonical-cost} to that same constant-after-$T$ gain. Thus
\begin{equation}\label{eq:actual-closed-cost}
 E\sup_{t\leq T}|M^{\rm actual}_t-M^{\rm actual}_0|
 +E\Var_{[0,T]}(A^{\rm actual})
 \leq30\bigl(32c+E|\Delta X_\tau|\bigr).
\end{equation}
Existence of the actual integral is supplied separately. This argument
transfers a quantitative bound to its components; uniqueness alone does not
provide that bound.
